% unit: u22
% sources: sources/086-cramer-s-rule-inverse-matrix-and-volume.txt, sources/039-lecture-20-cramer-s-rule-inverse-matrix-and-volume.txt, sources/191-mit18-06scf11-ses2-7sum-pdf-download-file.txt, sources/040-problem-solving-determinants-and-volume.txt, sources/192-mit18-06scf11-ses2-7prob-pdf-download-file.txt, sources/193-mit18-06scf11-ses2-7sol-pdf-download-file.txt
\section{Cramer's rule, inverse matrix and volume}\label{sec:u22}

The last two sections built the determinant from three properties and found two formulas
for it: the big formula with $n!$ terms and the cofactor expansion
(Section~\ref{sec:u21}). This section puts the determinant to work. The cofactor expansion
turns into an explicit formula for $A^{-1}$; applied to $\vb{b}$ that formula becomes
Cramer's rule for the solution of $A\vb{x}=\vb{b}$; and the absolute value $|\det A|$ turns
out to be a volume. Elimination remains the way to compute. What these formulas give
instead is \emph{algebra}: the answer written in terms of the entries, which shows how it
depends on them.

\subsection{The cofactor matrix gives \texorpdfstring{\(A^{-1}\)}{the inverse}}

Start from the $2\times 2$ inverse we already know:
\[
  \begin{bmatrix} a & b\\ c & d\end{bmatrix}^{-1}
  =\frac{1}{ad-bc}\begin{bmatrix} d & -b\\ -c & a\end{bmatrix}.
\]
We divide by the determinant, which is exactly what fails when $A$ is singular. The entries
of the matrix on the right are cofactors: $d$ is the cofactor of $a$ (strike out row~1 and
column~1), and $-b$ is the cofactor of $c$ (strike out row~2 and column~1, and attach the
sign $(-1)^{2+1}$). Notice that the cofactor of $c$, which sits in position $(2,1)$, appears
in position $(1,2)$. The cofactors are transposed.

\begin{definition}[Cofactor matrix]\label{def:u22-cofactor}
Let $A$ be $n\times n$ and let $M_{ij}$ be the $(n-1)\times(n-1)$ matrix left when row $i$
and column $j$ of $A$ are removed. The cofactor of $a_{ij}$ is
$C_{ij}=(-1)^{i+j}\det M_{ij}$, and the \emph{cofactor matrix} is $C=(C_{ij})$. Its
transpose $C\T$ is sometimes called the adjugate of $A$.
\end{definition}

Each cofactor is a sum of products of $n-1$ entries, while $\det A$ is a sum of products of
$n$ entries. So a formula $A^{-1}=C\T/\det A$ has the right ``degree'': entries of $A^{-1}$
behave like $1/(\text{entry of }A)$. For a general $3\times 3$ matrix with rows
$(a,b,c)$, $(d,e,f)$, $(g,h,i)$, the $(1,1)$ entry of $A^{-1}$ should be
$(ei-fh)/\det A$. The theorem says this pattern holds for every entry and every $n$.

\begin{theorem}[Inverse from cofactors]\label{thm:u22-inverse}
For every square matrix $A$,
\[
  A\,C\T=(\det A)\,I .
\]
Consequently, if $\det A\neq 0$ then $A^{-1}=\dfrac{1}{\det A}\,C\T$, that is,
$(A^{-1})_{ij}=C_{ji}/\det A$.
\end{theorem}

\begin{proof}
The $(i,k)$ entry of $A\,C\T$ is row $i$ of $A$ times column $k$ of $C\T$, which is row $k$
of $C$:
\[
  (A\,C\T)_{ik}=\sum_{j=1}^{n} a_{ij}\,C_{kj}.
\]
\emph{Diagonal entries} ($k=i$). The sum $\sum_j a_{ij}C_{ij}$ is the cofactor expansion of
$\det A$ along row $i$ (Section~\ref{sec:u21}). So every diagonal entry equals $\det A$.

\emph{Off-diagonal entries} ($k\neq i$). Let $A^{*}$ be the matrix $A$ with row $k$
replaced by a copy of row $i$. Then $A^{*}$ has two equal rows, so $\det A^{*}=0$. Expand
$\det A^{*}$ along its row $k$. The entries of that row are $a_{i1},\dots,a_{in}$, and the
cofactors of row $k$ of $A^{*}$ are the same as those of $A$, because they are computed after
deleting row $k$, the only row where $A$ and $A^{*}$ differ. Hence
\[
  0=\det A^{*}=\sum_{j} a_{ij}\,C_{kj}=(A\,C\T)_{ik}.
\]
So $A\,C\T$ has $\det A$ on its diagonal and zeros elsewhere. If $\det A\neq 0$, dividing
by $\det A$ gives $A\,(C\T/\det A)=I$, and a one-sided inverse of a square matrix is its
inverse.
\end{proof}

In the $2\times 2$ case the off-diagonal computation is visible: row~1 of $A$ is $(a,b)$,
the cofactors of row~2 are $(-b,a)$, and $a(-b)+b(a)=0$ is the determinant of the matrix
with two equal rows $\left[\begin{smallmatrix} a&b\\ a&b\end{smallmatrix}\right]$.

\begin{example}[A $3\times 3$ inverse by cofactors]\label{ex:u22-cofactors}
Let
\[
  A=\begin{bmatrix} 1&1&4\\ 1&2&2\\ 1&2&5\end{bmatrix}.
\]
The nine cofactors, each a signed $2\times 2$ determinant, are
\begin{align*}
C_{11}&=+(2\cdot5-2\cdot2)=6, & C_{12}&=-(1\cdot5-2\cdot1)=-3, & C_{13}&=+(1\cdot2-2\cdot1)=0,\\
C_{21}&=-(1\cdot5-4\cdot2)=3, & C_{22}&=+(1\cdot5-4\cdot1)=1, & C_{23}&=-(1\cdot2-1\cdot1)=-1,\\
C_{31}&=+(1\cdot2-4\cdot2)=-6, & C_{32}&=-(1\cdot2-4\cdot1)=2, & C_{33}&=+(1\cdot2-1\cdot1)=1.
\end{align*}
Multiplying $A$ by $C\T$,
\[
  A\,C\T=\begin{bmatrix} 1&1&4\\ 1&2&2\\ 1&2&5\end{bmatrix}
  \begin{bmatrix} 6&3&-6\\ -3&1&2\\ 0&-1&1\end{bmatrix}
  =\begin{bmatrix} 3&0&0\\ 0&3&0\\ 0&0&3\end{bmatrix}=3I,
\]
so $\det A=3$ and
\[
  A^{-1}=\frac13\begin{bmatrix} 6&3&-6\\ -3&1&2\\ 0&-1&1\end{bmatrix}
  =\begin{bmatrix} 2&1&-2\\ -1&\tfrac13&\tfrac23\\ 0&-\tfrac13&\tfrac13\end{bmatrix}.
\]
One more question the formula answers instantly: if the entry $a_{13}=4$ is changed to
$100$, the determinant does not change, because $\det A=a_{11}C_{11}+a_{12}C_{12}+a_{13}C_{13}$
and $C_{13}=0$. (The inverse does change: the cofactors in rows~2 and~3 involve $a_{13}$.)
\end{example}

\subsection{Cramer's rule}

If $A$ is invertible, the solution of $A\vb{x}=\vb{b}$ is
$\vb{x}=A^{-1}\vb{b}=C\T\vb{b}/\det A$. Cramer noticed that each component of $C\T\vb{b}$ is
itself a determinant.

\begin{theorem}[Cramer's rule]\label{thm:u22-cramer}
Let $A$ be $n\times n$ with $\det A\neq 0$, and let $B_j$ be the matrix $A$ with its column
$j$ replaced by $\vb{b}$. Then the solution of $A\vb{x}=\vb{b}$ is
\[
  x_j=\frac{\det B_j}{\det A},\qquad j=1,\dots,n.
\]
\end{theorem}

\begin{proof}
By Theorem~\ref{thm:u22-inverse}, $x_j=(C\T\vb{b})_j/\det A$ and
$(C\T\vb{b})_j=\sum_{i} C_{ij}\,b_i$. Now expand $\det B_j$ down its column $j$, which holds
$b_1,\dots,b_n$. The cofactor of position $(i,j)$ in $B_j$ is computed after deleting
column $j$, so it ignores $\vb{b}$ and equals the cofactor $C_{ij}$ of $A$. Therefore
$\det B_j=\sum_i b_i\,C_{ij}=(C\T\vb{b})_j$.
\end{proof}

\begin{example}[Cramer's rule on a $3\times 3$ system]\label{ex:u22-cramer}
Solve $A\vb{x}=\vb{b}$ with $A$ from Example~\ref{ex:u22-cofactors} and
$\vb{b}=(5,2,5)$. We know $\det A=3$. Replacing one column at a time,
\[
  \det B_1=\begin{vmatrix} 5&1&4\\ 2&2&2\\ 5&2&5\end{vmatrix}=6,\qquad
  \det B_2=\begin{vmatrix} 1&5&4\\ 1&2&2\\ 1&5&5\end{vmatrix}=-3,\qquad
  \det B_3=\begin{vmatrix} 1&1&5\\ 1&2&2\\ 1&2&5\end{vmatrix}=3 .
\]
For instance, expanding $\det B_1$ along row~1 gives
$5(10-4)-1(10-10)+4(4-10)=30-0-24=6$. Hence
$\vb{x}=(6/3,\,-3/3,\,3/3)=(2,-1,1)$. Check: $A(2,-1,1)=(2-1+4,\;2-2+2,\;2-2+5)=(5,2,5)$.
The same answer comes from $\vb{x}=C\T\vb{b}/3$:
$C\T\vb{b}=(30+6-30,\;-15+2+10,\;0-2+5)=(6,-3,3)$.
\end{example}

\begin{remark}[Formula versus algorithm]
Cramer's rule needs $n+1$ determinants of size $n$. Computed by elimination each costs about
$n^{3}/3$ operations, so the rule costs about $n^{4}/3$, against $n^{3}/3$ for solving the
system once by elimination; computed by the big formula it is hopeless. Software never uses
it. Its value is theoretical: it shows that each $x_j$ is a ratio of polynomials in the
entries of $A$ and $\vb{b}$, so the solution depends smoothly on the data wherever
$\det A\neq 0$, and it shows exactly how $\det A$ appears in every denominator.
\end{remark}

\begin{trap}
The inverse is $C\T/\det A$, not $C/\det A$. The cofactors of \emph{row}~$i$ of $A$ fill
\emph{column}~$i$ of $A^{-1}$. A quick test: for
$\left[\begin{smallmatrix} a&b\\ c&d\end{smallmatrix}\right]$ the $(1,2)$ entry of the
inverse must be $-b/(ad-bc)$, which is the cofactor of $c$, not of $b$. Also remember the
checkerboard signs $(-1)^{i+j}$; forgetting them is the other common error.
\end{trap}

\subsection{The determinant is a volume}

Take the rows of an $n\times n$ matrix $A$ as $n$ edges leaving the origin and complete
them to a box: the parallelogram ($n=2$), the parallelepiped ($n=3$), or its analogue in
$\R^{n}$, consisting of all points $t_1\vb{r}_1+\dots+t_n\vb{r}_n$ with $0\le t_i\le 1$.
The claim of this subsection is that the volume of that box is $|\det A|$.

We use only these geometric facts about volume. (V1) A rectangular box whose edges point
along the coordinate axes, with lengths $\ell_1,\dots,\ell_n$, has volume
$\ell_1\cdots\ell_n$. (V2) Listing the edges in another order gives the same box. (V3)
Volume is base times height: the box sits on the base spanned by $n-1$ of its edges, and its
height is the distance of the last edge's tip from that base. (V4) If the edges are dependent,
the box lies in a hyperplane and has volume $0$.

\begin{theorem}[Determinant and volume]\label{thm:u22-volume}
The box whose edges are the rows of $A$ has volume $|\det A|$. The same holds for the box
built on the columns, since $\det A\T=\det A$.
\end{theorem}

\begin{proof}
If $A$ is singular its rows are dependent, so the volume is $0=|\det A|$ by (V4).
Suppose $A$ is invertible. Elimination with row exchanges, followed by elimination upward,
reduces $A$ to a diagonal matrix $D=\operatorname{diag}(d_1,\dots,d_n)$ with nonzero
entries, using two kinds of steps (no row is ever rescaled).

\emph{Subtract $c$ times row $j$ from row $i$.} The determinant is unchanged (property~5 in
Section~\ref{sec:u20}). The box is sheared but its volume is unchanged: take as base the
face spanned by the other $n-1$ rows, which do not move. Row $i$ changes by a multiple of
row $j$, which lies in that base, so the tip of row $i$ moves parallel to the base and the
height stays the same. By (V3) the volume is unchanged.

\emph{Exchange two rows.} The determinant changes sign, so $|\det A|$ is unchanged, and by
(V2) the box is the same.

So both $|\det A|$ and the volume are unchanged at every step, and at the end both equal
$|d_1d_2\cdots d_n|$: the determinant of the diagonal matrix $D$ is the product of its
diagonal entries, and the box of $D$ is a rectangular box with edge lengths
$|d_1|,\dots,|d_n|$ along the axes (V1).
\end{proof}

The lecture argues differently but equivalently: the volume of the box satisfies the three
defining properties of $|\det|$. The unit cube ($A=I$) has volume $1$; exchanging two edges
does not change the box; stretching one edge by $t$ multiplies the volume by $|t|$; and
splitting one edge as a sum splits the volume (when the two pieces lie on the same side of the
base). The shear step in the proof is the geometric heart of the linearity property. The sign
of $\det A$ also means something: it records whether the edges form a right-handed or a
left-handed system, and exchanging two edges reverses the handedness.

A good test case is an orthogonal matrix $Q$. Its rows are perpendicular unit vectors, so its
box is a unit cube turned in space, with volume $1$. The determinant agrees: from
$Q\T Q=I$ and the product rule, $(\det Q\T)(\det Q)=(\det Q)^{2}=1$, so $\det Q=\pm 1$.

\begin{figure}[ht]\centering
\begin{tikzpicture}[scale=1.25,>=Stealth]
  % original parallelogram on r1=(3,1), r2=(1,2)
  \fill[blue!12] (0,0) -- (3,1) -- (4,3) -- (1,2) -- cycle;
  \fill[blue!30] (0,0) -- (3,1) -- (1,2) -- cycle;
  \draw[thick,blue!70!black] (0,0) -- (3,1) -- (4,3) -- (1,2) -- cycle;
  % sheared parallelogram on r1 and r2 + r1/2 = (2.5,2.5)
  \draw[thick,dashed,orange!80!black] (0,0) -- (3,1) -- (5.5,3.5) -- (2.5,2.5) -- cycle;
  % line through the tip of r2, parallel to r1
  \draw[gray,densely dotted] (-0.2,1.6) -- (6.1,3.7);
  \fill (0,0) circle (1.2pt) node[below left] {$O$};
  \draw[->,very thick] (0,0) -- (3,1) node[below right] {$\vb{r}_1=(a,b)=(3,1)$};
  \draw[->,very thick] (0,0) -- (1,2) node[left] {$\vb{r}_2=(c,d)=(1,2)$};
  \draw[->,very thick,orange!80!black] (0,0) -- (2.5,2.5) node[above left] {$\vb{r}_2+\tfrac12\vb{r}_1$};
  \node[blue!60!black] at (1.35,1.0) {\small triangle};
  \node[blue!60!black] at (3.0,2.0) {\small area $5$};
\end{tikzpicture}
\caption{The parallelogram with edges $\vb{r}_1=(3,1)$ and $\vb{r}_2=(1,2)$ has area
$|ad-bc|=|3\cdot 2-1\cdot 1|=5$, and the shaded triangle is half of it. Adding
$\tfrac12\vb{r}_1$ to $\vb{r}_2$ (a row operation) moves the tip along the dotted line
parallel to $\vb{r}_1$: the dashed parallelogram has the same base and height, hence the
same area, just as the determinant is unchanged.}
\label{fig:u22-area}
\end{figure}

\subsection{Areas and volumes from coordinates}

Theorem~\ref{thm:u22-volume} is a formula that needs no lengths, angles or heights, only
coordinates. In the plane, the parallelogram with edges $(a,b)$ and $(c,d)$ has area
$|ad-bc|$, and the triangle with vertices $(0,0)$, $(a,b)$, $(c,d)$ has half that area
(Figure~\ref{fig:u22-area}). For a triangle placed anywhere:

\begin{proposition}[Area of a triangle]\label{prop:u22-triangle}
The triangle with vertices $(x_1,y_1)$, $(x_2,y_2)$, $(x_3,y_3)$ has area
\[
  \frac12\left|\det\begin{bmatrix} x_1&y_1&1\\ x_2&y_2&1\\ x_3&y_3&1\end{bmatrix}\right| .
\]
\end{proposition}

\begin{proof}
Subtract row~1 from rows~2 and~3; the determinant does not change. The last column becomes
$(1,0,0)$, and expanding along it leaves
\[
  \det\begin{bmatrix} x_2-x_1 & y_2-y_1\\ x_3-x_1 & y_3-y_1\end{bmatrix}.
\]
The row operations translated the triangle so that its first vertex is at the origin; its
other vertices are now the two edge vectors, and half the parallelogram area is the triangle
area.
\end{proof}

For example, the triangle with vertices $(1,1)$, $(4,2)$, $(2,5)$ has edges $(3,1)$ and
$(1,4)$ from the first vertex, so its area is $\tfrac12|3\cdot4-1\cdot1|=\tfrac{11}{2}$. A
zero determinant means the three points are collinear.

\begin{example}[Volume of a tetrahedron]\label{ex:u22-tetra}
A tetrahedron $T$ has one vertex at the origin $O$ and the others at
$A_1=(2,2,-1)$, $A_2=(1,3,0)$ and $A_3=(-1,1,4)$. Find its volume. Then keep $A_1$ and
$A_2$, move the apex to $A_3'=(-201,-199,104)$, and find the new volume.

\emph{Relate $T$ to a box.} Let $P$ be the parallelepiped on the edges $OA_1$, $OA_2$,
$OA_3$. Take the triangle $OA_1A_2$ as the base of $T$ and the parallelogram on $OA_1$,
$OA_2$ as the base of $P$; the parallelogram has twice the area of the triangle, and both
solids have the same height $h$ (the distance from $A_3$ to that plane). So
\[
  \operatorname{vol}(T)=\tfrac13(\text{area of triangle})\,h
  =\tfrac16(\text{area of parallelogram})\,h=\tfrac16\operatorname{vol}(P).
\]
\emph{Compute.} By Theorem~\ref{thm:u22-volume},
\[
  \operatorname{vol}(P)=\left|\det\begin{bmatrix} 2&2&-1\\ 1&3&0\\ -1&1&4\end{bmatrix}\right|
  =\bigl|2(12-0)-2(4-0)+(-1)(1+3)\bigr|=|24-8-4|=12,
\]
so $\operatorname{vol}(T)=12/6=2$.

\emph{Move the apex.} Notice that $A_3'=A_3-100A_1$: indeed
$(-1,1,4)-100(2,2,-1)=(-201,-199,104)$. In the matrix this is the row operation ``subtract
$100$ times row~1 from row~3'', which leaves the determinant at $12$. So the long, needle-like
tetrahedron still has volume $2$. Geometrically, the apex moved parallel to the edge $OA_1$,
which lies in the base plane, so the height did not change.
\end{example}

\subsection{Exercises}

\begin{exercise}[Spherical coordinates]\label{exr:u22-spherical}
Spherical coordinates $(\rho,\phi,\theta)$ are defined by
\[
  x=\rho\sin\phi\cos\theta,\qquad y=\rho\sin\phi\sin\theta,\qquad z=\rho\cos\phi .
\]
Find the $3\times3$ matrix of partial derivatives of $(x,y,z)$ with respect to
$(\rho,\phi,\theta)$ and show that its determinant is $J=\rho^{2}\sin\phi$. Explain why
$\rho^{2}\sin\phi\,d\rho\,d\phi\,d\theta$ is the volume of a small coordinate box.
\end{exercise}
\begin{solution}
The rows are the partial derivatives of $x$, $y$, $z$:
\[
  J=\det\begin{bmatrix}
  \sin\phi\cos\theta & \rho\cos\phi\cos\theta & -\rho\sin\phi\sin\theta\\
  \sin\phi\sin\theta & \rho\cos\phi\sin\theta & \rho\sin\phi\cos\theta\\
  \cos\phi & -\rho\sin\phi & 0
  \end{bmatrix}.
\]
Expand along the bottom row (cofactor signs $+,-,+$; the last entry is $0$):
\begin{align*}
J&=\cos\phi\,\bigl(\rho^{2}\cos\phi\sin\phi\cos^{2}\theta+\rho^{2}\cos\phi\sin\phi\sin^{2}\theta\bigr)
 +\rho\sin\phi\,\bigl(\rho\sin^{2}\phi\cos^{2}\theta+\rho\sin^{2}\phi\sin^{2}\theta\bigr)\\
&=\rho^{2}\cos^{2}\phi\sin\phi+\rho^{2}\sin^{3}\phi
 =\rho^{2}\sin\phi\,(\cos^{2}\phi+\sin^{2}\phi)=\rho^{2}\sin\phi .
\end{align*}
Small changes $d\rho$, $d\phi$, $d\theta$ move the point by the three columns of this matrix
times $d\rho$, $d\phi$, $d\theta$. These three small vectors are the edges of the tiny
coordinate box, so by Theorem~\ref{thm:u22-volume} its volume is
$|J|\,d\rho\,d\phi\,d\theta=\rho^{2}\sin\phi\,d\rho\,d\phi\,d\theta$ (here $0\le\phi\le\pi$, so
$\sin\phi\ge0$).
\end{solution}

\begin{exercise}[Integer inverses]\label{exr:u22-integer}
Suppose $A$ has integer entries. Show that $A^{-1}$ has integer entries if $\det A=\pm1$,
and that conversely, if both $A$ and $A^{-1}$ have integer entries, then $\det A=\pm1$.
Find the inverse of $\left[\begin{smallmatrix} 2&3\\ 1&2\end{smallmatrix}\right]$.
\end{exercise}
\begin{solution}
Every cofactor is a sum of products of integers, so $C\T$ has integer entries; if
$\det A=\pm1$ then $A^{-1}=\pm C\T$ is an integer matrix. Conversely, $\det A$ and
$\det A^{-1}$ are integers (sums of products of integers) and their product is
$\det(AA^{-1})=\det I=1$, so both are $\pm1$. For the example, $\det A=4-3=1$ and
$A^{-1}=C\T=\left[\begin{smallmatrix} 2&-3\\ -1&2\end{smallmatrix}\right]$; indeed
$\left[\begin{smallmatrix} 2&3\\ 1&2\end{smallmatrix}\right]
\left[\begin{smallmatrix} 2&-3\\ -1&2\end{smallmatrix}\right]=I$.
\end{solution}

\begin{exercise}[Determinant of the cofactor matrix]\label{exr:u22-detC}
Show that $\det C=(\det A)^{n-1}$ for an invertible $n\times n$ matrix $A$, and check it on
Example~\ref{ex:u22-cofactors}.
\end{exercise}
\begin{solution}
Take determinants in $A\,C\T=(\det A)I$. The right side is a diagonal matrix with $n$ entries
equal to $\det A$, so its determinant is $(\det A)^{n}$. The left side gives
$\det A\cdot\det C\T=\det A\cdot\det C$. Dividing by $\det A\neq0$ gives
$\det C=(\det A)^{n-1}$. In Example~\ref{ex:u22-cofactors}, $n=3$ and $\det A=3$, so
$\det C$ should be $9$. Directly, expanding
$C=\left[\begin{smallmatrix} 6&-3&0\\ 3&1&-1\\ -6&2&1\end{smallmatrix}\right]$ along row~1:
$6(1+2)+3(3-6)+0=18-9=9$.
\end{solution}

\begin{exercise}[Cramer's rule and collinear points]\label{exr:u22-small}
(a) Solve $2x+y=3$, $x+3y=5$ by Cramer's rule.
(b) Use Proposition~\ref{prop:u22-triangle} to decide whether $(1,2)$, $(3,5)$, $(7,11)$
lie on one line.
\end{exercise}
\begin{solution}
(a) $\det A=2\cdot3-1\cdot1=5$, $\det B_1=\left|\begin{smallmatrix} 3&1\\ 5&3\end{smallmatrix}\right|=9-5=4$,
$\det B_2=\left|\begin{smallmatrix} 2&3\\ 1&5\end{smallmatrix}\right|=10-3=7$. So
$x=4/5$, $y=7/5$. Check: $8/5+7/5=3$ and $4/5+21/5=5$.

(b) Subtracting the first row as in the proof of Proposition~\ref{prop:u22-triangle} leaves
$\det\left[\begin{smallmatrix} 2&3\\ 6&9\end{smallmatrix}\right]=18-18=0$. The triangle
has area $0$, so the points are collinear: the edge $(6,9)$ is $3$ times the edge $(2,3)$.
\end{solution}
